
\documentclass[12pt,a4paper]{article}
\usepackage[utf8]{inputenc}
\usepackage[T1]{fontenc}
\usepackage{amsmath,amssymb,amsthm,amsfonts}
\usepackage{geometry}
\usepackage{graphicx}
\usepackage{booktabs}
\usepackage{hyperref}
\usepackage{enumitem}
\geometry{margin=1in}
\setlength{\parindent}{0pt}
\setlength{\parskip}{0.75em}

\title{Holonomic Relations: Core Formalism and Synthetic Validation}
\author{Adrian Lipa \and Ciel}
\date{\today}

\begin{document}
\maketitle

\begin{abstract}
This document presents a revised core formalism of the Holonomic Relations framework,
together with a synthetic validation protocol and a minimal Monte Carlo implementation.
The formalism is layered into geometry, interaction, and optional spectral structure.
The ideal Euler coherence constraint is distinguished from the measurable holonomy defect.
Synthetic experiments are included only as methodological validation of the internal dynamics;
they are not empirical confirmation.
\end{abstract}

\section{Structure of the formalism}

The framework is organised into three layers:
\begin{enumerate}[label=\textbf{L\arabic*.}]
\item \textbf{Geometry}: three fundamental phase modes on $\mathbb{CP}^2$.
\item \textbf{Interaction}: operational phases $(\gamma_S,\gamma_C,\gamma_Q,\gamma_T)$.
\item \textbf{Spectral (optional)}: auxiliary eigenvalues $\tau_i$ used as seeds.
\end{enumerate}

Only the geometric layer is strictly required for the core formalism.

\section{Geometric core}

\subsection{State space}

The relational manifold is taken to be
\[
\mathcal{M}_{\mathrm{rel}} = \mathbb{CP}^2,
\]
equipped with the Fubini--Study metric. In affine coordinates $z=(z_1,z_2)$:
\[
ds^2 = \frac{(1+\|z\|^2)\|dz\|^2 - |\bar z \cdot dz|^2}{(1+\|z\|^2)^2}.
\]

\subsection{Phase modes}

A geometric relational state is characterised by three fundamental phases
\[
\Gamma^{(3)}=(\gamma_1,\gamma_2,\gamma_3), \qquad \gamma_i \in S^1.
\]

\subsection{Ideal coherence and measurable defect}

The ideal Euler coherence constraint is
\[
\sum_{i=1}^3 e^{i\gamma_i}=0.
\]

In general, finite systems deviate from exact coherence. Define the holonomy defect
\[
\Delta_H := \sum_{i=1}^3 e^{i\gamma_i}
\]
and the observable holonomy density
\[
\mathcal{H}:=|\Delta_H|^2.
\]

Thus:
\begin{itemize}
\item $\mathcal{H}=0$ corresponds to exact coherence,
\item $\mathcal{H}>0$ measures deviation from the coherence manifold.
\end{itemize}

\section{Connection and white-thread amplitudes}

Assume a $U(1)$ connection $A$ with curvature
\[
F=dA=k\,\omega_{FS}, \qquad k\in\mathbb{Z}.
\]

For a path $\gamma_{ij}$, the associated phase transport is
\[
U_{ij} = \exp\!\left(i\int_{\gamma_{ij}} A\right).
\]

Define the white-thread amplitude by
\[
W_{ij}=\rho_{ij}U_{ij},
\qquad
\rho_{ij}=\cos\frac{d_{ij}}2\,e^{-\beta d_{ij}},
\]
where $d_{ij}$ is the Fubini--Study distance and $\beta>0$ is a damping parameter.

\section{Interaction layer}

We introduce the operational phases
\[
\gamma_S,\qquad \gamma_C,\qquad \gamma_Q,\qquad \gamma_T.
\]

A minimal projection into geometric modes is
\begin{align}
\gamma_1 &= \gamma_S - \gamma_T,\\
\gamma_2 &= \gamma_C - \gamma_T,\\
\gamma_3 &= \gamma_Q - \gamma_T.
\end{align}

This choice makes all operational phases relative to the truth phase $\gamma_T$, which acts as a gauge reference.

\section{Dynamics}

We define a Lagrangian
\[
\mathcal{L} = \mathcal{L}_{\mathrm{truth}} + \mathcal{L}_{\mathrm{coh}} + \mathcal{L}_{\mathrm{clarity}} - \mathcal{L}_{\mathrm{dist}}.
\]

The terms are
\begin{align}
\mathcal{L}_{\mathrm{truth}} &= -\frac12(\gamma_T-\gamma_T^{\mathrm{true}})^2,\\
\mathcal{L}_{\mathrm{coh}} &= -\frac12\sum_{i<j} C_{ij}^{(0)}(\gamma_i-\gamma_j)^2,\\
\mathcal{L}_{\mathrm{clarity}} &= -\frac12(\dot\gamma_S^2+\dot\gamma_C^2),\\
\mathcal{L}_{\mathrm{dist}} &= \sum_x \lambda_x \pi_x,
\end{align}
where
\[
\pi_x\in[0,1], \qquad x\in\{\mathrm{lie,omit,hallucinate,smooth}\},
\]
and $C_{ij}^{(0)}$ is evaluated in a frozen-coherence approximation at the previous step:
\[
C_{ij}^{(0)}=\frac{1+\cos(\gamma_i-\gamma_j)}{2}\Big|_{\mathrm{previous\ step}}.
\]

The overdamped dynamics is
\[
\dot\gamma_i = -\eta \frac{\partial V}{\partial \gamma_i},
\qquad
V=-\mathcal{L}_{\mathrm{truth}}-\mathcal{L}_{\mathrm{coh}}+\mathcal{L}_{\mathrm{dist}}.
\]

\section{Observables}

The two primary observables are:

\subsection{Holonomy density}
\[
\mathcal{H}=|\Delta_H|^2.
\]

\subsection{Semantic truth scalar}
\[
\Theta = 1-\frac{1}{|\mathcal{F}|}\sum_{f\in\mathcal{F}}
\bigl(\pi_{\mathrm{false}}(f)+\pi_{\mathrm{unmarked}}(f)\bigr),
\]
where $\mathcal{F}$ is the set of factual statements used in an audited interaction.

\section{Local stability statement}

Consider the overdamped dynamics in the distortion-free case. For sufficiently small step size and positive coherence couplings, the coherence manifold is a local attractor. The synthetic simulation below is intended as a numerical sanity check of this claim, not a proof of global convergence.

\section{Synthetic validation protocol}

Synthetic validation is used here only to test internal consistency.

\begin{enumerate}
\item Initialise random operational phases.
\item Project to geometric modes.
\item Evolve under coherence dynamics with additive Gaussian noise.
\item Introduce controlled distortion profiles.
\item Measure $\mathcal{H}(t)$, final $\mathcal{H}$, and $\Theta$.
\item Compare low-distortion and high-distortion regimes.
\end{enumerate}

\section{Results summary}

The accompanying code generates Monte Carlo trajectories and exports:
\begin{itemize}
\item mean holonomy decay,
\item final-state distribution,
\item a scatter plot of final $\mathcal{H}$ versus $\Theta$,
\item tabulated summary statistics.
\end{itemize}

These results are synthetic and therefore only support methodological consistency.

\section{Optional spectral seeds}

A convenient optional parametrisation is
\[
\tau_1=0.263,\qquad \tau_2=0.353,\qquad \tau_3=0.489.
\]
In this package, these values are treated as optional seeds rather than derivational necessities of the core formalism.

\section{Conclusion}

The package separates:
\begin{itemize}
\item the geometric core,
\item the operational interaction layer,
\item optional spectral structure,
\item and a synthetic validation module.
\end{itemize}

This separation makes the framework easier to audit, simulate, and modify.

\bibliographystyle{plain}
\bibliography{refs}

\end{document}
